\ifdefined\gkver\else\def\gkver{student}\fi
\documentclass[\gkver]{gaokaozhenti}
\begin{document}

\chapter{2005年天津理综}
%% paperType: 理综物理部分

\begin{enumerate}
\item
%% number: 14
%% paperName: 2005年天津理综第 14 题 6 分
%% typeId: 1
%% type: 单选题
%% chapter: 气体性质
%% point: 气体的状态参量
%% method: 无
%% score: 6
%% degree: 650
%% duplicateId: 0
%% body:
下列说法中正确的是\xzanswer{A}

\fourchoices[answer=A]
{一定质量的气体被压缩时，气体压强不一定增大}
{一定质量的气体温度不变压强增大时，其体积也增大}
{气体压强是由气体分子间的斥力产生的}
{在失重的情况下，密闭容器内的气体对器壁没有压强}

%% answer: A
%% memo:
\memoanswer{
本题原卷及材料中未提供详解，待补充。
}

\item
%% number: 15
%% paperName: 2005年天津理综第 15 题 6 分
%% typeId: 1
%% type: 单选题
%% chapter: 力物体的平衡
%% point: 多力平衡
%% method: 无
%% score: 6
%% degree: 650
%% duplicateId: 0
%% body:
如图所示，表面粗糙的固定斜面顶端安有滑轮，两物块 P、Q 用轻绳连接并跨过滑轮（不计滑轮的质量和摩擦），P 悬于空中，Q 放在斜面上，均处于静止状态。当用水平向左的恒力推 Q 时，P、Q 仍静止不动，则\xzanswer{D}

\onepicture[w=5.0cm]{figs/15.png}

\fourchoices[answer=D]
{Q 受到的摩擦力一定变小}
{Q 受到的摩擦力一定变大}
{轻绳上拉力一定变小}
{轻绳上拉力一定不变}

%% answer: D
%% memo:
\memoanswer{
本题原卷及材料中未提供详解，待补充。
}

\item
%% number: 16
%% paperName: 2005年天津理综第 16 题 6 分
%% typeId: 1
%% type: 单选题
%% chapter: 电磁感应
%% point: 切割磁感线电动势
%% method: 无
%% score: 6
%% degree: 650
%% duplicateId: 0
%% body:
将硬导线中间一段折成不封闭的正方形，每边长为 $l$，它在磁感应强度为 $B$、方向如图的匀强磁场中匀速转动，转速为 $n$，导线在 a、b 两处通过电刷与外电路连接，外电路有额定功率为 $P$ 的小灯泡并正常发光，电路中除灯泡外，其余部分的电阻不计，灯泡的电阻应为\xzanswer{B}

\onepicture[w=6.0cm]{figs/16.png}

\fourchoices[answer=B]
{$\frac{(2\pi l^{2}nB)^{2}}{P}$}
{$\frac{2(\pi l^{2}nB)^{2}}{P}$}
{$\frac{(l^{2}nB)^{2}}{2P}$}
{$\frac{(l^{2}nB)^{2}}{P}$}

%% answer: B
%% memo:
\memoanswer{
本题原卷及材料中未提供详解，待补充。
}

\item
%% number: 17
%% paperName: 2005年天津理综第 17 题 6 分
%% typeId: 1
%% type: 单选题
%% chapter: 光学
%% point: 光电效应
%% method: 无
%% score: 6
%% degree: 650
%% duplicateId: 0
%% body:
某光电管的阴极是用金属钾制成的，它的逸出功为 $2.21\UeV$，用波长为 $2.5\times10^{-7}\Um$ 的紫外线照射阴极，已知真空中光速为 $3.0\times10^{8}\Ums$，元电荷为 $1.6\times10^{-19}\UC$，普朗克常量为 $6.63\times10^{-34}\UJcdots$，求得钾的极限频率和该光电管发射的光电子的最大动能应分别是\xzanswer{B}

\fourchoices[answer=B]
{$5.3\times10^{14}\UHz$，$2.2\UJ$}
{$5.3\times10^{14}\UHz$，$4.4\times10^{-19}\UJ$}
{$3.3\times10^{33}\UHz$，$2.2\UJ$}
{$3.3\times10^{33}\UHz$，$4.4\times10^{-19}\UJ$}

%% answer: B
%% memo:
\memoanswer{
本题原卷及材料中未提供详解，待补充。
}

\item
%% number: 18
%% paperName: 2005年天津理综第 18 题 6 分
%% typeId: 1
%% type: 单选题
%% chapter: 电场
%% point: 带电粒子的轨迹类问题
%% method: 无
%% score: 6
%% degree: 650
%% duplicateId: 0
%% body:
一带电油滴在匀强电场 $E$ 中的运动轨迹如图中虚线所示，电场方向竖直向下。若不计空气阻力，则此带电油滴从 a 运动到 b 的过程中，能量变化情况为\xzanswer{C}

\onepicture[w=3.8cm]{figs/18.png}

\fourchoices[answer=C]
{动能减小}
{电势能增加}
{动能和电势能之和减小}
{重力势能和电势能之和增加}

%% answer: C
%% memo:
\memoanswer{
本题原卷及材料中未提供详解，待补充。
}

\item
%% number: 19
%% paperName: 2005年天津理综第 19 题 6 分
%% typeId: 1
%% type: 单选题
%% chapter: 机械波
%% point: 波的多解性
%% method: 无
%% score: 6
%% degree: 450
%% duplicateId: 0
%% body:
图中实线和虚线分别是 $x$ 轴上传播的一列简谐横波在 $t=0$ 和 $t=0.03\Us$ 时刻的波形图，$x=1.2\Um$ 处的质点在 $t=0.03\Us$ 时刻向 $y$ 轴正方向运动，则\xzanswer{A}

\onepicture[w=9.5cm]{figs/19.png}

\fourchoices[answer=A]
{该波的频率可能是 $125\UHz$}
{该波的波速可能是 $10\Ums$}
{$t=0$ 时 $x=1.4\Um$ 处质点的加速度方向沿 $y$ 轴正方向}
{各质点在 $0.03\Us$ 内随波迁移 $0.9\Um$}

%% answer: A
%% memo:
\memoanswer{
本题原卷及材料中未提供详解，待补充。
}

\item
%% number: 20
%% paperName: 2005年天津理综第 20 题 6 分
%% typeId: 1
%% type: 单选题
%% chapter: 电场
%% point: 带电粒子的直线运动
%% method: 无
%% score: 6
%% degree: 650
%% duplicateId: 0
%% body:
现用电子显微镜观测线度为 $d$ 的某生物大分子的结构。为满足测量要求，将显微镜工作时电子的德布罗意波长设定为 $\frac{d}{n}$，其中 $n>1$。已知普朗克常量 $h$、电子质量 $m$ 和电子电荷量 $e$，电子的初速度不计，则显微镜工作时电子的加速电压应为\xzanswer{D}

\fourchoices[answer=D]
{$\frac{n^{2}h^{2}}{med^{2}}$}
{$\left(\frac{md^{2}h^{2}}{n^{2}e^{3}}\right)^{\frac{1}{3}}$}
{$\frac{d^{2}h^{2}}{2men^{2}}$}
{$\frac{n^{2}h^{2}}{2med^{2}}$}

%% answer: D
%% memo:
\memoanswer{
本题原卷及材料中未提供详解，待补充。
}

\item
%% number: 21
%% paperName: 2005年天津理综第 21 题 6 分
%% typeId: 1
%% type: 单选题
%% chapter: 曲线运动
%% point: 人造地球卫星
%% method: 无
%% score: 6
%% degree: 450
%% duplicateId: 0
%% body:
土星周围有美丽壮观的“光环”，组成环的颗粒是大小不等、线度从 $1\Uum$ 到 $10\Um$ 的岩石、尘埃，类似于卫星，它们与土星中心的距离从 $7.3\times10^{4}\Ukm$ 延伸到 $1.4\times10^{5}\Ukm$。已知环的外缘颗粒绕土星做圆周运动的周期约为 $14\Uh$，引力常量为 $6.67\times10^{-11}\UNmqkgq$，则土星的质量约为（估算时不考虑环中颗粒间的相互作用）\xzanswer{D}

\fourchoices[answer=D]
{$9.0\times10^{16}\Ukg$}
{$6.4\times10^{17}\Ukg$}
{$9.0\times10^{25}\Ukg$}
{$6.4\times10^{26}\Ukg$}

%% answer: D
%% memo:
\memoanswer{
本题原卷及材料中未提供详解，待补充。
}

\item
%% number: 22
%% paperName: 2005年天津理综第 22 题 16 分
%% typeId: 4
%% type: 实验题
%% chapter: 光学
%% point: 双缝干涉测量波长
%% method: 无
%% score: 16
%% degree: 450
%% duplicateId: 0
%% body:
现有毛玻璃屏 A、双缝 B、白光光源 C、单缝 D 和透红光的滤光片 E 等光学元件，要把它们放在如图所示的光具座上组装成双缝干涉装置，用以测量红光的波长。

\begin{enumerate}
	\item
	将白光光源 C 放在光具座最左端，依次放置其他光学元件，由左至右，表示各光学元件的字母排列顺序应为 C、\tkanswer[1.5cm]{EDB}、A。

	\onepicture[w=9.0cm, align=right]{figs/22a.png}

	\item
	本实验的步骤有：
	\begin{steps}
		\item
		取下遮光筒左侧的元件，调节光源高度，使光束能直接沿遮光筒轴线把屏照亮；
		\item
		按合理顺序在光具座上放置各光学元件，并使各元件的中心位于遮光筒的轴线上；
		\item
		用米尺测量双缝到屏的距离；
		\item
		用测量头（其读数方法同螺旋测微器）测量数条亮纹间的距离。
	\end{steps}

	在操作步骤②时还应注意\tkanswer[1.6cm]{单缝和双缝间距 $5\Ucm\sim10\Ucm$}和\tkanswer[1.6cm]{使单缝与双缝相互平行}。

	\item
	将测量头的分划板中心刻线与某条亮纹中心对齐，将该亮纹定为第 1 条亮纹，此时手轮上的示数如图~\ref{2005天津22a} 所示。然后同方向转动测量头，使分划板中心刻线与第 6 条亮纹中心对齐，记下此时图~\ref{2005天津22b} 中手轮上的示数\tkanswer[2cm]{$13.870$}$\Umm$，求得相邻亮纹的间距 $\Delta x$ 为\tkanswer[2cm]{$2.310$}$\Umm$。

	\twopicture[labelA=2005天津22a, labelB=2005天津22b, align=right, h=2.6cm]
	{figs/22b.png}
	{figs/22c.png}

	\item
	已知双缝间距 $d$ 为 $2.0\times10^{-4}\Um$，测得双缝到屏的距离 $l$ 为 $0.700\Um$，由计算式 $\lambda=$\tkanswer[3cm]{$\frac{d}{l}\Delta x$}，求得所测红光波长为\tkanswer[2.5cm]{$6.6\times10^{2}$}$\Unm$。
\end{enumerate}

%% answer: 
\jdanswer{
\begin{enumerate}
	\item
	EDB

	\item
	单缝和双缝间距 $5\Ucm\sim10\Ucm$，使单缝与双缝相互平行

	\item
	13.870，2.310

	\item
	$\frac{d}{l}\Delta x$，$6.6\times10^{2}$
\end{enumerate}
}

%% memo:
\memoanswer{
本题原卷及材料中未提供详解，待补充。
}

\item
%% number: 23
%% paperName: 2005年天津理综第 23 题 16 分
%% typeId: 6
%% type: 计算题
%% chapter: 电磁感应
%% point: 电磁感应能量分析
%% method: 无
%% score: 16
%% degree: 450
%% duplicateId: 0
%% body:
图中 MN 和 PQ 为竖直方向的两平行长直金属导轨，间距 $l$ 为 $0.40\Um$，电阻不计。导轨所在平面与磁感应强度 $B$ 为 $0.50\UT$ 的匀强磁场垂直。质量 $m$ 为 $6.0\times10^{-3}\Ukg$、电阻为 $1.0\UO$ 的金属杆 ab 始终垂直于导轨，并与其保持光滑接触。导轨两端分别接有滑动变阻器和阻值为 $3.0\UO$ 的电阻 $R_{1}$。当杆 ab 达到稳定状态时以速率 $v$ 匀速下滑，整个电路消耗的电功率 $P$ 为 $0.27\UW$，重力加速度取 $10\Umsq$，试求速率 $v$ 和滑动变阻器接入电路部分的阻值 $R_{2}$。

\onepicture[h=4.2cm, align=right]{figs/23.png}

%% answer: 
\jdanswer{
\begin{enumerate}
	\item
	$v=4.5\Ums$

	\item
	$R_{2}=6.0\UO$
\end{enumerate}
}

%% memo:
\memoanswer{
由能量守恒定律得：$mgv=P$ ①

代入数据得：$v=4.5\Ums$ ②

又：$E=BLv$ ③

设电阻 $R_{a}$ 与 $R_{b}$ 的并联电阻为 $R_{\text{外}}$，ab 棒的电阻为 $r$，有

$\frac{1}{R_{\text{外}}}=\frac{1}{R_{a}}+\frac{1}{R_{b}}$ ④

$I=E/(R_{\text{外}}+r)$ ⑤

$P=IE$ ⑥

代入数据得：$R_{2}=6.0\UO$ ⑦
}

\item
%% number: 24
%% paperName: 2005年天津理综第 24 题 18 分
%% typeId: 6
%% type: 计算题
%% chapter: 运动和力
%% point: 动量定理
%% method: 无
%% score: 18
%% degree: 450
%% duplicateId: 0
%% body:
如图所示，质量 $m_{A}$ 为 $4.0\Ukg$ 的木板 A 放在水平面 C 上，木板与水平面间的动摩擦因数 $\mu$ 为 0.24，木板右端放着质量 $m_{B}$ 为 $1.0\Ukg$ 的小物块 B（视为质点），它们均处于静止状态。木板突然受到水平向右的 $12\UNs$ 的瞬时冲量 $I$ 作用开始运动，当小物块滑离木板时，木板的动能 $E_{kB}$ 为 $8.0\UJ$，小物块的动能为 $0.50\UJ$，重力加速度取 $10\Umsq$，求：

\begin{enumerate}
	\item
	瞬时冲量作用结束时木板的速度 $v_{0}$；
	\item
	木板的长度 $L$。
\end{enumerate}

\onepicture[w=6.5cm, align=right]{figs/24.png}

%% answer: 
\jdanswer{
\begin{enumerate}
	\item
	$v_{0}=3.0\Ums$

	\item
	$L=0.50\Um$
\end{enumerate}
}

%% memo:
\memoanswer{
（1）设水平向右为正方向，有：$I=m_{A}v_{0}$ ①

代入数据得：$v=3.0\Ums$ ②

（2）设 A 对 B、B 对 A、C 对 A 的滑动摩擦力的大小分别为 $F_{AB}$、$F_{BA}$、$F_{CA}$，B 在 A 上滑行的时间为 $t$，B 离开 A 时 A 和 B 的速度分别为 $v_{A}$ 和 $v_{B}$，有

$-\left(F_{BA}+F_{CA}\right)t=mv_{A}-m_{A}v_{0}$ ③

$F_{AB}t=m_{B}v_{B}$ ④

其中 $F_{AB}=F_{BA}$

$F_{CA}=\mu(m_{A}+m_{C})g$ ⑤

设 A、B 相对于 C 的位移大小分别为 $s_{A}$ 和 $s_{B}$，有

$-\left(F_{BA}+F_{CA}\right)s_{A}=mv_{A}^{2}/2-m_{A}v_{0}^{2}/2$ ⑥

$F_{AB}s_{B}=E_{KB}$ ⑦

动量和动能之间的关系为：$m_{A}v_{A}=\sqrt{2m_{A}E_{KA}}$ ⑧

$m_{A}v_{A}=\sqrt{2m_{A}E_{KA}}$ ⑨

木板 A 的长度：$L=s_{A}-s_{B}$ ⑩

代入数据得：$L=0.50\Um$
}

\item
%% number: 25
%% paperName: 2005年天津理综第 25 题 22 分
%% typeId: 6
%% type: 计算题
%% chapter: 磁场
%% point: 洛伦兹力
%% method: 无
%% score: 22
%% degree: 450
%% duplicateId: 0
%% body:
正电子发射计算机断层（PET）是分子水平上的人体功能显像的国际领先技术，它为临床诊断和治疗提供全新的手段。

\begin{enumerate}
	\item
	PET 在心脏疾病诊疗中，需要使用放射正电子的同位素氮 13 示踪剂。氮 13 是由小型回旋加速器输出的高速质子轰击氧 16 获得的，反应中同时还产生另一个粒子，试写出该核反应方程。
	\item
	PET 所用回旋加速器示意如图，其中置于高真空中的金属 D 形盒的半径为 $R$，两盒间距为 $d$，在左侧 D 形盒圆心处放有粒子源 S，匀强磁场的磁感应强度为 $B$，方向如图所示。质子质量为 $m$，电荷量为 $q$。设质子从粒子源 S 进入加速电场时的初速度不计，质子在加速器中运动的总时间为 $t$（其中已略去了质子在加速电场中的运动时间），质子在电场中的加速次数与回旋半周的次数相同，加速质子时的电压大小可视为不变。求此加速器所需的高频电源频率 $f$ 和加速电压 $U$。
	\item
	试推证当 $R\gg d$ 时，质子在电场中加速的总时间相对于在 D 形盒中回旋的时间可忽略不计（质子在电场中运动时，不考虑磁场的影响）。
\end{enumerate}

\onepicture[h=5.0cm, align=right]{figs/25.png}

%% answer: 
\jdanswer{
\begin{enumerate}
	\item
	\ce{^{16}_{8}O + ^{1}_{1}H -> ^{13}_{7}N + ^{4}_{2}He}

	\item
	$f=\frac{qB}{2\pi m}$，$U=\frac{\pi BR^{2}}{2t}$

	\item
	$\frac{t_{1}}{t_{2}}=\frac{2d}{\pi R}\ll1$
\end{enumerate}
}

%% memo:
\memoanswer{
（1）核反应方程为：$\ce{^{16}_{8}O + ^{1}_{1}H -> ^{13}_{7}N + ^{4}_{2}He}$ ①

（2）设质子加速后最大速度为 $v$，由牛顿第二定律得：$qvB=mv^{2}/R$ ②

质子的回旋周期为：$T=2\pi R/v=2\pi m/qB$ ③

高频电源的频率为：$f=1/T=\frac{qB}{2\pi m}$ ④

质子加速后的最大动能为：$E_{k}=mv^{2}/2$ ⑤

设质子在电场中加速的次数为 $n$，则：$E_{k}=nqU$ ⑥

又 $t=nT/2$ ⑦

可解得：$U=\frac{\pi BR^{2}}{2t}$ ⑧

（3）在电场中加速的总时间为：$t_{1}=2nd/v$ ⑨

在 D 形盒中回旋的总时间为 $t_{2}=n\pi R/v$ ⑩

故：$\frac{t_{1}}{t_{2}}=\frac{2d}{\pi R}\ll1$，即当 $R\gg d$ 时，可以忽略不计。
}

\end{enumerate}

\end{document}
